% Model Checker: a report for Zenodo.
% Build: pdflatex, bibtex, pdflatex, pdflatex (main file modelchecker.tex),
% or: latexmk -pdf modelchecker.tex
\documentclass[10pt]{article}
\usepackage[a4paper,margin=1.8cm,top=1.6cm]{geometry}
\usepackage{multicol}
\usepackage[T1]{fontenc}
\usepackage{lmodern}
\usepackage{microtype}
\usepackage{booktabs}
\usepackage{tabularx}
\usepackage{url}
\usepackage[hidelinks]{hyperref}
\usepackage[numbers,sort&compress]{natbib}
\usepackage[compact]{titlesec}
\titleformat{\section}{\normalfont\large\bfseries}{\thesection}{0.6em}{}
\titlespacing*{\section}{0pt}{0.9em}{0.2em}
\setlength{\bibsep}{0pt plus 0.3ex}
\setlength{\parskip}{0.25em}
\setlength{\parindent}{0pt}

\title{\vspace{-2.5em}Model Checker: Static Analysis of Kinetic SBML Models}
% Inside \author there can be no blank lines: use \\ for line breaks.
\author{Herbert M. Sauro$^{1,2}$\\[0.5ex]
  \small $^{1}$Department of Bioengineering and $^{2}$eScience Institute, University of Washington, Seattle, WA, 98195-5061, USA
}
\date{\vspace{-1ex}\small Version 1.0.0, November 2026. DOI: \href{https://doi.org/10.5281/zenodo.24265079}{10.5281/zenodo.24265079}}

\begin{document}
\maketitle
\vspace{-2em}

\begin{abstract}
\noindent
Model Checker examines kinetic models written in SBML for errors that a schema or consistency validator does
not detect but that cause a simulation to fail or to give misleading results. Such errors include rate laws that cannot be
evaluated, reactions that can drive a species negative, rate laws inconsistent with their reaction's
stoichiometry, inconsistent units, thermodynamically inconsistent reversible reactions, and invalid or
misplaced annotations. The current version performs only static checks, that is it examines the model's equations and structure without
simulating it. Three techniques are used. Interval arithmetic evaluates each rate law over the whole range of
values its variables can take, which shows whether the rate can become undefined (for example by division by
zero) or can continue to consume a species that has reached zero. The stoichiometry matrix gives the
conservation laws, the species that can never be produced, and the cycles of reversible reactions whose
equilibrium constants are tested for thermodynamic consistency. Dimensional analysis compares the units of the
terms in each expression. Where possible, a finding includes a specific example, such as values of the
species at which a rate law divides by zero, so that the problem can be confirmed directly. The software is written in Rust and runs either as a command-line program or
in the web browser, compiled to WebAssembly; in the browser the analysis runs locally and the model is not sent
to a server. The web version also shows the model as Antimony text beside the findings, so that the model can be
corrected and checked again; edits are merged back into the original SBML.

\noindent\textbf{Availability:} \url{https://sys-bio.github.io/ModelChecker/} (web application);\\
\url{https://github.com/sys-bio/ModelChecker} (source, MIT licence).
\end{abstract}

\section{Background}
SBML \cite{hucka2003sbml} is the common exchange format for kinetic models in systems biology. libSBML
\cite{bornstein2008libsbml} validates a document against the specification, checking syntax, references
and, where units are declared, unit consistency. A model can pass these checks and still be unusable: a
Michaelis--Menten denominator that is zero at the initial state, a degradation reaction written as a
constant rate that continues to consume a species after it reaches zero, or a reversible rate law whose
equilibrium constant depends on an effector. Such errors are usually found only when a simulation fails, and may then be
difficult to trace to their source. MEMOTE
\cite{lieven2020memote} provides a test suite of this kind for genome-scale constraint-based models; Model
Checker addresses kinetic models, in which the rate laws themselves are the main source of error.

\section{Checks}
Version 1.0.0 has 80 finding codes in twelve groups (Table~\ref{tab:checks}). Each finding has a severity:
\emph{error} (a simulation will fail, or the model is ill-defined), \emph{warning} (the model is probably not
what was intended) or \emph{info} (worth knowing; often harmless). Function definitions and assignment rules
are substituted into the expressions that use them; species are non-negative unless they start negative.

\begin{table}[ht]
\centering\footnotesize
\begin{tabularx}{\textwidth}{@{}lrX@{}}
\toprule
Group & Codes & Examples \\
\midrule
Rate laws (RL) & 13 & rate cannot be evaluated at the initial state or at a reachable state; species can become negative; zero-order consumption; reactant missing from the rate law; mass-action powers differ from stoichiometry \\
Rules and events (RU) & 9 & cyclic assignment or initial-assignment rules; variable set more than once; constant set by a rule; event trigger true at $t=0$ or never able to fire \\
Network (NW) & 7 & conservation laws; species that can never be produced and the reactions that therefore never run; disconnected subnetworks \\
Thermodynamics (TD) & 2 & Haldane relation violated; Wegscheider condition violated around a cycle of reversible reactions \\
Units (UN) & 8 & rate law not in extent per time; sum or comparison of different dimensions or scales; missing compartment volume factor \\
Values (VA) & 6 & parameter value inconsistent with its role (e.g.\ negative $K_m$); reaction with zero rate; rate constants spanning ten or more orders of magnitude \\
SBO terms (SO) & 8 & invalid or obsolete term; term from the wrong branch; kinetic law labelled mass action or Hill that does not have that form \\
Annotations (MD) & 7 & missing database links or model metadata; identifier invalid for its database; qualifier inappropriate for the element \\
Symbols (SY) & 7 & undefined, unused or shadowed symbols; symbol with no value \\
Other (SB, FT, AN) & 13 & structural problems; features many simulators do not support (fast reactions, delays, algebraic rules, SBML packages); Antimony errors \\
\bottomrule
\end{tabularx}
\vspace{-0.8em}
\caption{Groups of checks. The full list, with severities, is in the repository README.}
\label{tab:checks}
\end{table}
\vspace{-1.2em}

\section{Methods}
\textbf{Evaluability.} Each rate law and rule is evaluated by interval arithmetic \cite{moore2009interval}
over the states the model could reach, which shows in one evaluation whether a division by zero, logarithm of
zero or similar fault is possible. If so, a concrete point is sought. A fault at the initial values is an error, a fault at some other reachable state a warning, and a
fault that interval arithmetic cannot rule out but for which no point is found is reported for information
only, because interval bounds overestimate.

\textbf{Positivity.} For a species $S$ with net stoichiometry $n$ in a reaction of rate $v$, the model
preserves $S \ge 0$ only if $n\,v \ge 0$ whenever $S = 0$. This is tested from the rate law alone; the
\texttt{reversible} attribute is not used. A counterexample is first sought with the other variables between
zero and their initial values; one found only outside that range is reported for information, since it
commonly reflects a part exceeding its total (as in $E_T - ES$).

\textbf{Network structure.} Conservation laws are computed from the left null space of the stoichiometry
matrix \cite{sauro2004conservation} and
expressed with few negative coefficients ($S + ES + P$). Species that can never become non-zero are found by a fixed-point
calculation related to the empty-siphon analysis of Petri nets \cite{angeli2007petri}. Starting with all
species at zero, a species is released when some reaction can produce it, where whether a reaction can run,
and in which direction, is decided from its rate law by interval arithmetic rather than from its reactant list.

\textbf{Thermodynamics.} For a reversible rate law depending on both reactants and products, the mass-action
ratio at which the rate is zero is located numerically at several random states. A consistent rate law gives
a single value, its equilibrium constant. Around each cycle of reversible reactions (from the null space of the stoichiometry matrix with
boundary species) the product of equilibrium constants must be 1 within 1\%.

\textbf{Units and annotations.} Units are inferred from declared units as SI exponents with a scale factor;
a parameter without units makes any expression containing it unknown, so unit findings are raised only where
the model declares enough to support them. Database links \cite{lenovere2005miriam} are checked against a
built-in copy of the identifiers.org registry \cite{wimalaratne2018identifiers}, for a known namespace and
an identifier of the correct form, and \texttt{bqbiol:is} links for the kind of entity they name. SBO terms
\cite{courtot2011sbo} are checked against a built-in copy of the ontology, both for the branch the element
permits and, for kinetic laws, against the form of the rate law.

\section{Implementation and testing}
The software is written in Rust as five crates: an SBML reader and writer (Levels 2 and 3), an Antimony
\cite{smith2009antimony} parser, the checks (without input or output, so that they compile to WebAssembly),
a command-line program for files or directories, and the web page. The web page shows the model as editable Antimony model
beside the findings. The original SBML remains the reference copy. An edit is merged by comparing the edited
text with the text as first generated, so content the text does not represent (notes, other tools'
annotations, layouts) is retained. On the curated BioModels collection \cite{malik2020biomodels}, unedited text reproduced the
original document exactly for all 1,075 models, and single edits of a value, rate law, SBO term, unit or
database link changed only the edited item in every model tested. The repository has 141 automated tests.

\section{Limitations}
Findings should be read as evidence for
inspection, not as proof. Interval arithmetic overestimates ranges; the thermodynamic tests sample a finite
number of states; unit checks are only as complete as the units the model declares. Reachability is approximated, not computed, so
a reported counterexample state may not be reachable from the initial conditions. SBML Level 1 is not read, and of hierarchical (\texttt{comp}) models only the top-level model
is checked. The rate of false positives on a large corpus has not yet been measured systematically.

{\footnotesize
% TODO (author): review and rewrite this section in your own words, as for the PathwayAnnotator report.
\textbf{Use of AI.} Claude Code (Anthropic, Opus 5.5) was used to assist in writing the software and in preparing this report,
including writing and debugging code, writing tests, and managing commits to the GitHub repository. The test
suite and the Rust linter Clippy were run on the code before each commit. The author reviewed, verified and
edited all generated content and takes full responsibility for the integrity and accuracy of the final
publication. Lessons for AI-assisted development as reported for Iridium~\cite{iridium2026} were employed
throughout.

\textbf{Acknowledgements.} This work was supported by NIH Biomedical Imaging and Bioengineering award P41 EB023912 through HMS at the
Center for Reproducible Biomedical Modeling (\url{https://reproduciblebiomodels.org}).
}

{\footnotesize
\begin{multicols}{2}
\bibliographystyle{unsrtnat}
\bibliography{references}
\end{multicols}
}

\end{document}
